\slajdid{08-s022}
\begin{frame}[label=08-s022]
Zbog osobine regenerabilnosti logičkih kola smemo dozvoliti smetnje koje su i veće nego što smo prethodno definisali. Praktično smetnja, ako se pojavljuje samo u jednoj tački, sme da poremeti nivo praktično do napona $V_M$ pošto će se naponski nivoi regenerisati. U tom slučaju možemo da definišemo margine šuma za jednostruke izvore šuma (SS – single source) kao
\[NM_{LSS}=V_M-V_{OL}\]
\[NM_{HSS}=V_{OH}-V_M\]
\par\bigskip
Međutim, u realnom radu ne možemo očekivati da će se smetnje pojavljivati samo u jednoj tački. Odnosno ostaje pitanje kako odrediti napone $V_{IL}$ i $V_{IH}$ ako dozvolimo da se smetnje mogu pojavljivati u više tačaka u kolu.
\end{frame}
